\documentclass[10pt,a4paper]{amsart}
\usepackage{amsmath,amssymb,amsthm,graphicx}
% Retained minimal preamble for the Stacks Project chapter more-morphisms.tex (More on Morphisms).
% Verbatim excerpt of the chapter preamble: Xy-pic setup (native preamble lines 14-18) and the native
% theorem declarations with the mathematical macros (native preamble lines 154-217). Copied mechanically;
% nothing added, removed or redefined. The wrapper class and the amsmath/amssymb/amsthm/graphicx base
% packages are supplied by the assembler, so the class/externaldocument part of the native preamble is omitted.
% For commutative diagrams we use Xy-pic
\usepackage[all]{xy}

% We use 2cell for 2-commutative diagrams.
\xyoption{2cell}
% Theorem environments.
%
\theoremstyle{plain}
\newtheorem{theorem}[subsection]{Theorem}
\newtheorem{proposition}[subsection]{Proposition}
\newtheorem{lemma}[subsection]{Lemma}

\theoremstyle{definition}
\newtheorem{definition}[subsection]{Definition}
\newtheorem{example}[subsection]{Example}
\newtheorem{exercise}[subsection]{Exercise}
\newtheorem{situation}[subsection]{Situation}

\theoremstyle{remark}
\newtheorem{remark}[subsection]{Remark}
\newtheorem{remarks}[subsection]{Remarks}

\numberwithin{equation}{subsection}

% Macros
%
\def\lim{\mathop{\mathrm{lim}}\nolimits}
\def\colim{\mathop{\mathrm{colim}}\nolimits}
\def\Spec{\mathop{\mathrm{Spec}}}
\def\Hom{\mathop{\mathrm{Hom}}\nolimits}
\def\Ext{\mathop{\mathrm{Ext}}\nolimits}
\def\SheafHom{\mathop{\mathcal{H}\!\mathit{om}}\nolimits}
\def\SheafExt{\mathop{\mathcal{E}\!\mathit{xt}}\nolimits}
\def\Sch{\mathit{Sch}}
\def\Mor{\mathop{\mathrm{Mor}}\nolimits}
\def\Ob{\mathop{\mathrm{Ob}}\nolimits}
\def\Sh{\mathop{\mathit{Sh}}\nolimits}
\def\NL{\mathop{N\!L}\nolimits}
\def\CH{\mathop{\mathrm{CH}}\nolimits}
\def\proetale{{pro\text{-}\acute{e}tale}}
\def\etale{{\acute{e}tale}}
\def\QCoh{\mathit{QCoh}}
\def\Ker{\mathop{\mathrm{Ker}}}
\def\Im{\mathop{\mathrm{Im}}}
\def\Coker{\mathop{\mathrm{Coker}}}
\def\Coim{\mathop{\mathrm{Coim}}}

% Boxtimes
%
\DeclareMathSymbol{\boxtimes}{\mathbin}{AMSa}{"02}

%
% Macros for moduli stacks/spaces
%
\def\QCohstack{\mathcal{QC}\!\mathit{oh}}
\def\Cohstack{\mathcal{C}\!\mathit{oh}}
\def\Spacesstack{\mathcal{S}\!\mathit{paces}}
\def\Quotfunctor{\mathrm{Quot}}
\def\Hilbfunctor{\mathrm{Hilb}}
\def\Curvesstack{\mathcal{C}\!\mathit{urves}}
\def\Polarizedstack{\mathcal{P}\!\mathit{olarized}}
\def\Complexesstack{\mathcal{C}\!\mathit{omplexes}}
% \Pic is the operator that assigns to X its picard group, usage \Pic(X)
% \Picardstack_{X/B} denotes the Picard stack of X over B
% \Picardfunctor_{X/B} denotes the Picard functor of X over B
\def\Pic{\mathop{\mathrm{Pic}}\nolimits}
\def\Picardstack{\mathcal{P}\!\mathit{ic}}
\def\Picardfunctor{\mathrm{Pic}}
\def\Deformationcategory{\mathcal{D}\!\mathit{ef}}
\usepackage{cleveref}
\crefname{originallemma}{original lemma}{original lemmas}
\Crefname{originallemma}{Original lemma}{Original lemmas}
\makeatletter
\newlabel{lemma-smoothness-stratification-at-generic-point}{{0H43}{source}{}{sourcepointer-1}{}}
\newlabel{lemma-smoothness-stratification-at-generic-point@cref}{{[originallemma][0][]0H43}{[1][0][]source}{}{}{}}
\makeatother
\crefname{originalreference}{original reference}{original references}
\Crefname{originalreference}{Original reference}{Original references}
\makeatletter
\newlabel{limits-lemma-descend-finite-presentation}{{01ZM}{source}{}{sourcepointer-2}{}}
\newlabel{limits-lemma-descend-finite-presentation@cref}{{[originalreference][0][]01ZM}{[1][0][]source}{}{}{}}
\makeatother
\crefname{originalreference}{original reference}{original references}
\Crefname{originalreference}{Original reference}{Original references}
\makeatletter
\newlabel{limits-proposition-approximate}{{01ZA}{source}{}{sourcepointer-3}{}}
\newlabel{limits-proposition-approximate@cref}{{[originalreference][0][]01ZA}{[1][0][]source}{}{}{}}
\makeatother
\title{Stratification of a finitely presented morphism by smoothness loci}
\author{The Stacks Project Authors}
\date{stacks-project:more-morphisms.tex:local-snapshot-20261011; GNU FDL 1.2 or later}
\begin{document}
\maketitle
\noindent\textit{Editorial note.} Counted span is the author's own proof environment, source lines 16011-16052, immediately following the statement at 15986-16009.\par\medskip
\noindent\textit{Editorial note.} Two cross-chapter pointers (Limits) plus one same-chapter lemma are recorded in assembly.external\_references.\par\medskip
\noindent\textit{Editorial note.} The statement's descending chain notation with its conditions is retained verbatim in the statement span.\par\medskip
\noindent\textit{Editorial note.} Steps the author leaves to the reader are flagged here for the reviewer: source line(s) 16024 carry an "omitted"/"follows easily"/"immediate from the definitions" phrase. Those sentences are retained verbatim; the displayed argument around them stays inside the counted span.\par\medskip
\noindent\textit{Allowed standard prerequisites (editorial).}\begin{itemize}
\item Smoothness stratification at a generic point. chapter lemma lemma-smoothness-stratification-at-generic-point, cited in the proof. Supplies the local step: after shrinking, the morphism becomes smooth over a nonempty open of the source.
\item Approximation of a finitely presented morphism by Noetherian data. Limits, proposition limits-proposition-approximate and lemma limits-lemma-descend-finite-presentation, cited in the proof. Lets the Noetherian fppf-local statements be applied after approximating the situation.
\item Noetherian induction on the source. standard argument used to remove each successive smoothness locus. Iterates the local step to produce the descending chain of closed subschemes.
\end{itemize}\medskip
\ifdefined\hypertarget\hypertarget{sourcepointer-1}{}\fi
\noindent\textit{Source reference.} Original lemma 0H43: sources/stacks-more-morphisms/pinned-original-source.tex, source lines 15886--15886. Navigation pointer only: the counted proof cites this statement by \textbackslash{}ref as a standing input; the deduction the counted proof has to make is carried out inside the counted span, and no part of that deduction is delegated to this pointer.\par\medskip
\ifdefined\hypertarget\hypertarget{sourcepointer-2}{}\fi
\noindent\textit{Source reference.} Original reference 01ZM: focused-seeds/geometry/S9-morph-smoothness-stratification/ref-tags.txt, source lines 4--4. Navigation pointer only: the counted proof cites this statement by \textbackslash{}ref as a standing input; the deduction the counted proof has to make is carried out inside the counted span, and no part of that deduction is delegated to this pointer.\par\medskip
\ifdefined\hypertarget\hypertarget{sourcepointer-3}{}\fi
\noindent\textit{Source reference.} Original reference 01ZA: focused-seeds/geometry/S9-morph-smoothness-stratification/ref-tags.txt, source lines 5--5. Navigation pointer only: the counted proof cites this statement by \textbackslash{}ref as a standing input; the deduction the counted proof has to make is carried out inside the counted span, and no part of that deduction is delegated to this pointer.\par\medskip
\begin{lemma}
\label{lemma-smoothness-stratification}
Let $f : X \to S$ be a morphism of finite presentation between quasi-compact
and quasi-separated schemes. Then there exists a $t \geq 0$ and closed
subschemes
$$
S \supset S_0 \supset S_1 \supset \ldots \supset S_t = \emptyset
$$
such that
\begin{enumerate}
\item $S_i \to S$ is defined by a finite type ideal sheaf,
\item $S_0 \subset S$ is a thickening,
\item for each $i$ there exists a surjective finite locally free
morphism $T_i \to S_i \setminus S_{i + 1}$,
\item for each $i$ there exists a $t_i \geq 0$ and closed subschemes
$$
X_i = X \times_S T_i \supset Z_{i, 0}
\supset Z_{i, 1} \supset \ldots \supset Z_{i, t_i} = \emptyset
$$
such that $Z_{i, j} \to X_i$ is defined by a finite type ideal sheaf,
$Z_{i, 0} \subset X_i$ is a thickening, and such that the morphism
$Z_{i, j} \setminus Z_{i, j + 1} \to T_i$ is smooth.
\end{enumerate}
\end{lemma}

\begin{proof}
We can find a cartesian diagram
$$
\xymatrix{
X \ar[d] \ar[r] & X_0 \ar[d] \\
S \ar[r] & S_0
}
$$
such that $X_0$ and $S_0$ are of finite type over $\mathbf{Z}$. See
Limits, Proposition \ref{limits-proposition-approximate} and
Lemma \ref{limits-lemma-descend-finite-presentation}.
Thus we may assume $X$ and $S$ are of finite type over $\mathbf{Z}$.
Namely, a solution of the problem posed by the lemma for $X_0 \to S_0$
will base change to a solution over $S$; details omitted.

\medskip\noindent
Assume $X$ and $S$ are of finite type over $\mathbf{Z}$.
In this case every quasi-coherent ideal is of finite type, hence
we do not have to check the condition that $S_i$ is cut out
by a finite type ideal. Set $S_0 = S_{red}$ equal to the reduction of $S$.
Let $\eta \in S_0$ be a generic point of an irreducible component.
By Lemma \ref{lemma-smoothness-stratification-at-generic-point}
we can find an open subscheme $U \subset S_0$, 
a surjective, universally injective, finite locally free
morphism $V \to U$, a $t_0 \geq 0$ and closed subschemes
$$
X \times_S V \supset Z_{0, 0} \supset Z_{0, 1} \supset \ldots \supset
Z_{0, t_0} = \emptyset
$$
such that $Z_{0, i} \to X \times_S V$ is defined by a finite type ideal sheaf,
$Z_{0, 0} \subset X \times_S V$ is a thickening, and such that the morphism
$Z_{0, i} \setminus Z_{0, i + 1} \to V$ is smooth.
Then we let $S_1 \subset S_0$ be the reduced induced subscheme structure
on $S_0 \setminus U$. By Noetherian induction on the underlying
topological space of $S$, we may assume that the lemma holds for
$X \times_S S_1 \to S_1$. This produces $t \geq 1$ and
$$
S_1 = S_1 \supset S_2 \supset \ldots \supset S_t = \emptyset
$$
and $t_i$ and $Z_{i, j}$ as in the statement of the lemma.
This proves the lemma.
\end{proof}

\end{document}
